% Generated draft: requires mathematical and source review.
\documentclass[11pt]{article}
\usepackage[margin=25mm]{geometry}
\usepackage{fontspec}
\setmainfont{Latin Modern Roman}
\usepackage{amsmath,amssymb,mathtools,unicode-math}
\setmathfont{Latin Modern Math}
\usepackage{enumitem,needspace,xurl,hyperref}
\setlength{\parindent}{0pt}
\setlength{\parskip}{5pt}
\setlength{\emergencystretch}{3em}
\newcommand{\N}{\mathbb{N}}
\newcommand{\Z}{\mathbb{Z}}
\newcommand{\Q}{\mathbb{Q}}
\newcommand{\R}{\mathbb{R}}
\newcommand{\C}{\mathbb{C}}
\newcommand{\F}{\mathbb{F}}
\newcommand{\A}{\mathbb{A}}
\newcommand{\PP}{\mathbb{P}}
\newcommand{\E}{\mathbb{E}}
\newcommand{\eps}{\varepsilon}
\newcommand{\dd}{\,\mathrm{d}}
\begin{document}
\section*{UnsolvedMath ID: 30006648}
\textbf{Source problem code:} catalogue entry
\subsection{Definitions and mathematical statement}
Let \(k\) be an algebraically closed field of prime characteristic \(p>0\), and let \(G\) be a simple simply connected algebraic group over \(k\), with root system \(\Phi\) and Coxeter number \(h\). The supplied statement asks for two assertions about Donkin's Tilting Module Conjecture, but does not state that conjecture or specify its Frobenius-kernel scope. The following gives an explicit conventional interpretation; its agreement with the intended source formulation requires human review.

Choose a maximal torus and a positive system of roots. Write \(X^+\) for the dominant integral weights, \(\rho\) for half the sum of the positive roots, and \(w_0\) for the longest Weyl-group element. For an integer \(r\geq 1\), let \(G_r=\ker(F^r:G\longrightarrow G^{(r)})\) be the \(r\)-th Frobenius kernel, and put
\[
X_r=\{\lambda\in X^+:0\leq\langle\lambda,\alpha^\vee\rangle<p^r\text{ for every simple root }\alpha\}.
\]
For \(\lambda\in X_r\), let \(L_G(\lambda)\) be the simple rational \(G\)-module of highest weight \(\lambda\), and let \(Q_r(\lambda)\) be the projective cover of \(L_G(\lambda)|_{G_r}\) in the category of finite-dimensional \(G_r\)-modules. For a dominant weight \(\mu\), let \(T_G(\mu)\) denote the indecomposable rational tilting \(G\)-module of highest weight \(\mu\); a tilting module admits both a Weyl filtration and a dual Weyl filtration. Define the assertion
\[
\mathcal D_r(G):\quad
T_G\bigl(2(p^r-1)\rho+w_0\lambda\bigr)|_{G_r}
\cong Q_r(\lambda)
\qquad\text{for every }\lambda\in X_r.
\]
For a fixed root-system type \(\Phi\) and a prime \(p\), write \(\mathcal D(\Phi,p)\) for the assertion that \(\mathcal D_r(G)\) holds for every integer \(r\geq1\), every algebraically closed field \(k\) of characteristic \(p\), and every simple simply connected \(G/k\) of type \(\Phi\).

\textbf{Requested assertions, conditional on this interpretation.} For every irreducible root-system type \(\Phi\), prove
\[
\bigl(\forall\text{ primes }p,\ \mathcal D(\Phi,p)\bigr)
\quad\Longleftrightarrow\quad
\bigl(\Phi\text{ has type }A_n\text{ for some }n\geq1\text{ or type }B_2\bigr),
\]
and prove
\[
\forall\text{ primes }p,\qquad p>h\ \Longrightarrow\ \mathcal D(\Phi,p).
\]
The first assertion concerns validity for all prime characteristics at a fixed type; it does not assert failure at every prime for other types. The second assertion uses the strict inequality \(p>h\) and makes no assertion at \(p=h\).

\textbf{Unresolved source qualification.} The supplied original excerpt describes a broad problem of lifting modular representations from Lie algebras to algebraic groups and names a contribution on Donkin's conjectures; it contains neither the two requested assertions nor the displayed module isomorphism. The reconstructed statement names Donkin's Tilting Module Conjecture without specifying whether it means \(r=1\), all \(r\geq1\), or a formulation retaining a torus action. The definition of \(\mathcal D(\Phi,p)\) above is therefore an explicit candidate interpretation, not a source-verified identification. Both requested assertions remain conjectural targets in this draft.
\subsection{Short English statement}
The supplied target says that Donkin's Tilting Module Conjecture should hold in every prime characteristic exactly for types \(A_n\) and \(B_2\), and should hold in every type when \(p>h\). An explicit formulation using projective modules for Frobenius kernels is given conditionally, because the supplied excerpt does not specify the precise version of the conjecture. The supplied open status and reconstruction claim require independent review.
\subsection{Sources}
\begin{itemize}
\item Supplied local UnsolvedMath registry record: \texttt{data/unsolvedmath/problems.json}, permanent ID \(30006648\), title ``Revised Donkin Tilting-Module Conjectures,'' exact source problem code \texttt{catalogue entry}; collection: Oberwolfach Reports Open Problems. The record supplies the year \(2026\), labels the problem ``open,'' and labels its reconstructed statement \texttt{corrected\_verified}, with the claim that it was checked against the official report. These are supplied registry assertions, not independent verification performed for this draft.
\item \emph{Cohomology of Finite Groups: Interactions and Applications} (2026), Oberwolfach Reports, pp.\ 541--610, DOI \texttt{10.4171/owr/2026/9}. Supplied originating-report link: \url{https://doi.org/10.4171/owr/2026/9}. The supplied original excerpt identifies p.\ 549 and the contribution ``Cohomology for Algebraic Groups and Donkin's Conjectures'' by Daniel K. Nakano, reporting joint work with Christopher P. Bendel, Cornelius Pillen, and Paul Sobaje. The record does not separately identify a proposer.
\item Bendel--Nakano--Pillen--Sobaje, \emph{Donkin's Tilting Module Conjecture I}. Supplied literature link: \url{https://arxiv.org/abs/2103.14164}.
\item Bendel--Nakano--Pillen--Sobaje, \emph{Donkin's Tilting Module Conjecture II}. Supplied literature link: \url{https://arxiv.org/abs/2107.11615}.
\item Supplied literature-check date: \texttt{2026-08-22}. This is supplied dated curation, not an independently verified literature assessment. That curation labels the problem open, describes the extracted statement as a broad lifting problem rather than a theorem, and reports substantial case-by-case progress and counterexamples to nearby formulations, with no general resolution. It expressly requests independent verification of time-sensitive status and cited claims. No source was visited in preparing this draft.
\end{itemize}
\end{document}
